%! Author = sriadityavedantam
%! Date = 3/30/26

\section{Symmetric and Alternating Groups}

\begin{theorem*}{
    Every element in $S_n$ can be written as a product of disjoint cycles.
}
\end{theorem*}

\[
    (123)(4567)
\]
is disjoint.
\[
    (123)(426)
\]
is not disjoint.

\begin{theorem}{
    Every $\omega\in S_n$ can be written as a product of transpositions.
}
    For example,
    \begin{align*}
        \sigma &= (12\dots n) \\
        &= (1n)(1(n-1))\dots(13)(12).
    \end{align*}
\end{theorem}

\begin{theorem}{
    Suppose $\sigma\in S_n$ with
    \begin{align*}
        \sigma &= \tau_{1}\dots \tau_{k}, \qquad \tau_{i}\text{ transpositions}, \\
        &= \lambda_{1}\dots \lambda_{r}, \qquad \lambda_{i}\text{ transpositions}.
    \end{align*}
    Then
    \[
        k\equiv r\mod 2.
    \]
}
    Suppose $\tau=(ab)$ is a transposition.
    Then
    \[
        \tau^{-1}=\tau.
    \]
    Suppose
    \[
        \sigma=\tau_{1}\dots \tau_{k}
    \]
    is a product of transpositions.
    Then
    \[
        \sigma^{-1}=\tau_{k}\dots \tau_{1}.
    \]
    First reduce to the case $\sigma=e$.
    Suppose
    \[
        \sigma=\tau_{1}\dots \tau_{k}=\lambda_{1}\dots \lambda_{r}.
    \]
    Then
    \[
        1=\sigma\sigma^{-1}=\tau_{1}\dots \tau_{k}\lambda_{r}\dots \lambda_{1}.
    \]
    So the identity is written as a product of $k+r$ transpositions.
    Thus it is enough to prove that whenever
    \[
        e=\tau_{1}\dots \tau_{m},
    \]
    the integer $m$ is even.
    Assume for contradiction that
    \[
        e=\tau_{1}\dots \tau_{m}
    \]
    with $m$ odd.
    Choose such an expression with the fewest number of transpositions possible.
    Some $a\in\{1,\dots,n\}$ must appear in one of the transpositions.
    Among all such expressions, choose one where the leftmost occurrence of $a$ is as far to the right as possible.
    Now examine the transpositions immediately around that occurrence.
    Using the relations among transpositions, we can commute disjoint transpositions and simplify adjacent ones.
    If two equal transpositions appear, they cancel.
    If a transposition involving $a$ is followed by another transposition sharing exactly one element, then their product can be rewritten in a way that moves the occurrence of $a$ farther to the right or reduces the total number of transpositions.
    Either way, this contradicts the minimal choice of the expression.
    Hence the identity cannot be written as a product of an odd number of transpositions.
    Therefore $k+r$ is even, so
    \[
        k\equiv r\mod 2.
    \]
\end{theorem}

\begin{definition}[Sign]
    The sign of $\sigma$ is $1$ if $\sigma$ can be written as a product of an even number of transpositions, and $-1$ otherwise.
    The sign of $\sigma$ is denoted by $\sgn(\sigma)$ or $\epsilon(\sigma)$.
    The previous theorem says $\sgn(\sigma)$ is well-defined.
\end{definition}

\begin{definition}[Even Symmetric Group]
    \[
        A_n \coloneqq \{\sigma\in S_n : \sgn(\sigma)=1\}.
    \]
\end{definition}

\begin{theorem}{
    $A_n$ is a subgroup of $S_n$.
}
    First, $e\in A_n$ since
    \[
        \sgn(e)=1.
    \]
    Now let $\sigma,\tau\in A_n$.
    Then both $\sigma$ and $\tau$ can be written as products of an even number of transpositions.
    Therefore $\sigma\tau$ can also be written as a product of an even number of transpositions.
    Hence
    \[
        \sgn(\sigma\tau)=1,
    \]
    so $\sigma\tau\in A_n$.
    Also, if $\sigma\in A_n$, then $\sigma$ is a product of an even number of transpositions.
    Since
    \[
        \sigma^{-1}
    \]
    is obtained by reversing that product, $\sigma^{-1}$ is also a product of an even number of transpositions.
    Thus
    \[
        \sgn(\sigma^{-1})=1,
    \]
    so $\sigma^{-1}\in A_n$.
    Therefore $A_n\leq S_n$.
\end{theorem}